\documentclass[11pt,a4paper]{article}
\usepackage[margin=1in]{geometry}
\usepackage{booktabs}
\usepackage{amsmath,amssymb}
\usepackage{longtable}
\usepackage[dvipsnames]{xcolor}
\usepackage{hyperref}
\hypersetup{colorlinks=true,linkcolor=blue!50!black,urlcolor=blue!50!black}
\usepackage{parskip}

\title{Filtering the bracketing phase\\
\large A rectangle benchmark for candidate eigenvalue searches in \texttt{chladni}}
\author{}
\date{2 September 2026}

% \pr is a built-in relation symbol; use \pr for the precision parameter
\newcommand{\pr}{\ensuremath{\mathrm{prec}}}
\newcommand{\lam}{\ensuremath{\lambda}}
\newcommand{\sig}{\ensuremath{\sigma}}

\begin{document}
\maketitle

\begin{abstract}
\noindent
An eigenvalue search in \texttt{chladni} decomposes into three phases: \emph{bracketing}
(produce intervals each believed to hold one eigenvalue), \emph{refinement} (locate the minimum
inside a bracket), and \emph{multiplicity estimation}. This report tests the first phase alone,
on rectangles, where the spectrum is known in closed form and clustering and multiplicity are
under exact control.

Bracketing is the phase worth testing first. Every search failure measured during the design
work was a bracketing failure, and bracketing is where essentially all the evaluations go ---
refinement is a handful of them and multiplicity is three. Testing it in isolation also
sidesteps a defect in the joint scorer, described in \S\ref{sec:whybracket}, which makes
recall-with-multiplicity unable to distinguish ``split the pair'' from ``guessed~2''.

The purpose is to \textbf{filter}: to identify methods not worth developing further. Results are
therefore stated as capability claims within an explicit evaluation budget, never as performance
scores.
\end{abstract}

\tableofcontents
\newpage

\section{What is being tested, and why bracketing alone}
\label{sec:whybracket}

\subsection{The three phases}

Most candidate searches decompose the same way: a bracketing pass produces candidate intervals,
a local refinement locates the minimum inside each, and a multiplicity estimator reads the
tension spectrum at the located point. This decomposition was verified rather than assumed. In
all five searches implemented to date, the bracketing pass provably terminates before the first
call to the refinement routine, and no refined value re-enters subdivision; none is entangled.

The three phases are also very unequal. All five searches dispatch to one of only two
multiplicity estimators, and only one of those is live --- the ratio-gap plus local-minimum test
(C1)+(C2). So multiplicity is a \emph{shared} component: a defect there affects every candidate
identically, and the filter cannot discriminate between candidates on that axis. Refinement is
similarly settled: Brent on $\sig_1^2$ is measured machine-exact on every domain tried, at every
regularization tolerance from \texttt{1e-7} to \texttt{1e-16}. That leaves bracketing as the
phase where the candidates actually differ.

\subsection{The joint scorer cannot measure what we need}

The end-to-end scorer matches found eigenvalues against truth at a fixed relative tolerance of
\texttt{1e-6}, expanding a reported $(\lam, \mu)$ into $\mu$ copies of $\lam$. Below
$d/\lam < \texttt{1e-6}$, a method that never separates a pair but claims multiplicity~2 has both
copies match two \emph{distinct} truth entries, scoring recall $1.000$, zero spurious, and
\texttt{perfect=True} --- while the honest report of one root of multiplicity~1 scores $0.500$
\emph{and} takes a multiplicity penalty. Measured:

\begin{center}
\begin{tabular}{rrll}
\toprule
$d$ & $d/\lam$ & reports one root, claims $\mu=2$ & reports one root, honest $\mu=1$ \\
\midrule
5.9e-04 & 1.2e-05 & recall 0.000, spurious 2 & recall 0.000, spurious 1 \\
5.9e-05 & 1.2e-06 & recall 1.000, \textbf{perfect} & recall 0.500 \\
5.9e-07 & 1.2e-08 & recall 1.000, \textbf{perfect} & recall 0.500, multX 1 \\
\bottomrule
\end{tabular}
\end{center}

\noindent
The metric does not merely go blind, it rewards the wrong behaviour, and it does so at
$\pr\!\cdot\!\lam \approx \texttt{4.9e-5}$, essentially the boundary at which the basis stops
being able to resolve the pair at all. This is the same error the reference layer had already
found and removed from its own \texttt{collapse} function, three orders tighter, with the
comment that a tolerance ``does not describe the spectrum, it imposes a policy on it''.

A bracket set carries no multiplicities, so the exploit does not exist and the joint scorer is
not on the path for anything in this report.

\section{The acceptance rule}

\subsection{A lattice, not a bijection}

Let $S(\pr)$ denote the true spectrum with levels merged whenever a group's \emph{diameter}
is at most $\pr\!\cdot\!\lam$. A bracket set is acceptable when it lies in the lattice

\begin{center}
\begin{tabular}{ll}
\textbf{fine end} & one bracket per true eigenvalue \\
\textbf{coarse end} & one bracket per level of $S(\pr)$ \\
\end{tabular}
\end{center}

\noindent
and anything between the two is acceptable. In particular an \emph{accidental} split of a
within-$\pr$ cluster is not penalised: downstream, refinement finds both members, the
collapse rule sees two roots inside the merge tolerance, and the reported answer is two levels
of multiplicity~1 --- which is legal, because merging is permitted and never required.
Penalising the bracketer for it would contradict the multiplicity convention.

Bounding the bracket count rather than fixing it is also what defeats the two degenerate
strategies. Returning $[(a,b)]$ scores coverage $1.0$ and zero spurious while being useless;
returning ten thousand slivers has perfect purity and absurd cost. Neither survives a rule that
constrains the count from both sides.

\subsection{Metrics}

\begin{description}
\item[Coverage] --- hard. Every true eigenvalue lies in some bracket. A miss is unrecoverable
downstream.
\item[Purity] --- hard, and the metric with teeth. No bracket may contain eigenvalues spanning
more than $\pr\!\cdot\!\lam$. If it does, refinement returns one of them and the other is
lost.
\item[Fineness] --- free within the lattice. Reported against both bounds.
\item[Spurious] --- a count, not a failure: cost, plus the risk of a phantom root from a bracket
that brushes a well without containing its minimum.
\item[Cost] --- evaluations. Tension and eigenvector calls are counted in one counter, since an
eigenvector call costs 100--107\,\% of a tension call and a method that leans on them should
not get them free.
\end{description}

\noindent
A method \textbf{passes} a case when coverage is $1.0$ and no bracket is impure.

\subsection{Committing}

A bracketing method must \emph{commit}: the returned intervals are those it believes hold an
eigenvalue, not a subdivision of $[a,b]$. Without this requirement a partitioning method gets
coverage for free. Measured on one case with four true eigenvalues in $[10, 88]$: the swing's
cells tile the interval, total width $78.0$ of $78.0$, and score coverage $4/4$ by
construction, while uniform windows cover $5.9$ of $78.0$ and have to earn it.

\section{The instrument}

\subsection{Why rectangles}

The search layer is basis-agnostic by construction --- it sees only the tension function and the
interval. So a method can only respond to the \emph{shape} of the tension curve, and that shape
has four parameters: the floor $\varepsilon$, the wall slope $c$, the plateau, and the eigenvalue
spacing. Rectangles control all four, and add three things no other domain offers: closed-form
truth (no table-precision ceiling), exact control of clustering and multiplicity, and low cost.

On the unit square the eigenvalues are $\lam/\pi^2 = m^2 + n^2$, so multiplicity is a
number-theoretic property of the level: $2$ at $N=5$, $3$ at $N=50$, $4$ at $N=65$, $6$ at
$N=325$. Perturbing to $\mathrm{rect}(1+\delta, 1)$ splits a multiplicity-$m$ level into $m$
eigenvalues whose spacing scales with $\delta$ and whose internal gaps are unequal --- for
$N=65$ they are in the ratio $1 : 2.2 : 1$, and that ratio is invariant as $\delta$ sweeps
decades. The cluster \emph{shape} is fixed while its \emph{scale} varies, so a method cannot
pass by exploiting one particular spacing.

\subsection{The basis, and a correction}
\label{sec:basis}

The pencil must resolve the tension to below the finest $\pr$ under test, or the experiment
measures the objective rather than the algorithm. Checking that requirement found that
\textbf{corner-adapted Fourier--Bessel is the wrong basis for a rectangle}, by four orders.
Unit square, truncation effectively off, $\sig_\mu$ at each level (i.e.\ counting multiplicity):

\begin{center}
\begin{tabular}{lrrrrr}
\toprule
basis & cols & $N{=}2\,(\mu{=}1)$ & $N{=}5\,(\mu{=}2)$ & $N{=}50\,(\mu{=}3)$ & $N{=}65\,(\mu{=}4)$ \\
\midrule
FB $n{=}80$   &  80 & 1.68e-15 & 2.60e-14 & 1.42e-11 & 3.17e-12 \\
FB $n{=}160$  & 160 & 1.56e-15 & 2.51e-14 & 1.41e-11 & 3.15e-12 \\
PlaneWave 100 & 200 & 4.57e-16 & 5.67e-16 & 1.67e-15 & 1.66e-15 \\
PlaneWave 160 & 320 & 4.86e-16 & 6.02e-16 & 1.20e-15 & 1.45e-15 \\
\bottomrule
\end{tabular}
\end{center}

\noindent
The reason is that a right angle has $\pi/\omega = 2$, an integer, so the eigenfunction is
\emph{analytic} there and no fractional expansion is called for. Corner-FB is the right tool for
a singular corner and the wrong one for a domain with none; that pure Fourier--Bessel wins on
the L-shape does not transfer.

Three candidate explanations for the Fourier--Bessel floor were tested and rejected before the
basis itself was suspected. It is not the expansion order (flat from 80 to 300 columns), not the
collocation (flat from $n_b = 400$ to $6400$ and $n_i = 200$ to $3200$), and not the rank
truncation --- sweeping the threshold from its default \texttt{1.5e-12} down to zero moves the
floor by under 1\,\% while the retained rank rises from 70 to 80. It is also not a property of
$\lam$: the mode $(7,1)$ gives the same $3.5$--$3.7\mathrm{e}{-12}$ on squares of side 1, 2, 3
and 5, i.e.\ at $\lam$ from 493 down to 19.7, so it is set by $k\!\cdot\!\mathrm{diam}$, the mode
index, and rescaling the domain cannot buy precision.

With sized plane waves every level used in this report resolves to \texttt{1.2e-15} or better,
multiplicity~6 included.

\subsection{The controlling principle: fix the pencil, vary only \texorpdfstring{$\pr$}{prec}}

\textbf{The pencil is built once per case at the best resolution available and is identical
across the whole $\pr$ sweep.} $\pr$ enters only as an input to the bracketing algorithm,
telling it when to stop splitting.

This is the difference between testing bracketing and re-discovering the limits of an
under-resolved tension curve. Three separate basis-construction errors arose during this work
--- a rank frozen over a wide window, too few Fourier--Bessel terms on a corner domain, and
corner-FB on a rectangle --- and every one silently degraded the objective by orders while
looking healthy. With the pencil fixed, any difference across $\pr$ is attributable to the
algorithm by construction. A hard precondition asserts, before any case is scored, that
$\sig_\mu$ is at or below \texttt{1e-14} at every true eigenvalue counting multiplicity, and
that the finest $\pr$ in the sweep clears the pencil's own limit.

\subsection{Truth}

Window completeness comes from the closed-form spectrum function, which is complete by
construction rather than by truncation. The perturbed levels take their values from the mode
formula directly, because the spectrum function snaps near-degenerate values together at a
relative tolerance of \texttt{1e-13} --- measured, it holds the $5\pi^2$ pair distinct down to
$d = \texttt{5.9e-12}$ and merges it at \texttt{6.0e-13}. No case in this report approaches that
floor.

\section{The methods}

Each is exposed as \texttt{brackets(tensions, a, b, *, prec, modes)}, extracted from the
corresponding search rather than reimplemented, and reusing that search's own helpers. The five
searches themselves are untouched.

\begin{description}
\item[\texttt{uniform\_grid}] One grid, discrete minima of $\sig_1$, bracket from the
neighbouring samples. Deliberately not $\pr$-aware beyond a cap: a uniform grid cannot reach a
spacing of $\pr\!\cdot\!\lam$ at any usable cost, so its brackets are grid-spacing wide
whatever is asked. That is the honest behaviour of the method.

\item[\texttt{lappy}] The original recursive bisection. Its recursion is gated on $\sig_2$'s
\emph{value} against a threshold, which is a property of the grid rather than of the pair: a
grid of spacing $h$ sees about $c h/2$, so the gate engages when $h \le 2\,\mathrm{ttol}/c$
whatever the separation.

\item[\texttt{lappy\_repaired}] The same, with the gate removed: recurse into \emph{every}
discrete minimum, unconditionally, down to a floor. The floor here is $\pr\!\cdot\!\lam$
rather than a fixed fraction of the interval --- the fixed fraction was previously found to be
load-bearing and unexplained, losing a different eigenvalue at each of three values.

\item[\texttt{lipschitz}] Piyavskii--Shubert branch and bound. Its completeness certificate is
conditional on the Lipschitz constant $L$, and $L$ estimated from a pilot has been measured
\emph{unsound}: on one domain a 20-point pilot with $\gamma = 2$ underestimated the true slope,
ratio $0.672$, so the certificate can be false. It is included to measure that, not because it
is expected to win. Brackets are read off the sample cache, not off the candidate intervals ---
the candidate list in the original is dead code, left behind when refinement was rewritten to
work from the samples.

\item[\texttt{swing\_cells}] Cells cut at peaks of the angle between successive minimising
eigenvectors. Those peaks fall at branch \emph{crossings}, which lie between eigenvalues, so
they partition the interval into cells holding one eigenvalue each. Each crossing is
\textbf{bisected} rather than taken at the peak sample: detection localises a crossing only to
the sampling interval, and when the pair separation is below the stride the cell boundary can
land outside the pair. Bisection costs $\log_2(h/d)$ eigenvector evaluations and converges
regardless of how coarse the stride was.

\item[\texttt{swing\_direct}] The same cells, each then narrowed by a Lipschitz-free
\textsc{direct} subdivision before the bracket is returned. The ablation against
\texttt{swing\_cells} is what attributes the composed method's behaviour to one mechanism or
the other.
\end{description}

\subsection{Controls}

Two adversarial controls, without which the metrics are unvalidated:
\texttt{single\_bracket} returns $[(a,b)]$ and must fail purity on every case with more than one
level; \texttt{fine\_grid} returns thousands of equal slivers and must fail on spurious count
and cost. If either scores well, the metric is broken rather than the control.

\section{Results}

\texttt{orc} is the cheapest of the tested settings that reached coverage $1.0$ with no impure bracket; \texttt{aut} is the single setting fixed across the whole suite before any per-case result was seen. \textbf{--} means no setting succeeded within the budget of 1500 evaluations. Only the \texttt{aut} column may be aggregated: \texttt{orc} is a best-of-$K$ statistic and is biased upward in $K$.

\subsection{Aggregate: autonomous passes}

Across the 27 swept cases (the multiplicity-6 stress case excluded), counting only the autonomous setting:

\begin{center}
\begin{tabular}{lrrrr}
\toprule
method & $\pr=10^{-4}$ & $10^{-8}$ & $10^{-11}$ & $10^{-14}$ \\
\midrule
\texttt{uniform\_grid} & 16/27 & 9/27 & 7/27 & 5/27 \\
\texttt{lappy} & 17/27 & 10/27 & 8/27 & 6/27 \\
\texttt{lappy\_repaired} & 12/27 & 5/27 & 5/27 & 5/27 \\
\texttt{lipschitz} & 22/27 & 14/27 & 12/27 & 10/27 \\
\texttt{swing\_cells} & 21/27 & 13/27 & 11/27 & 8/27 \\
\texttt{swing\_direct} & 22/27 & 14/27 & 10/27 & 8/27 \\
\midrule
\texttt{single\_bracket} & 5/27 & 3/27 & 3/27 & 3/27 \\
\texttt{fine\_grid} & 0/27 & 0/27 & 0/27 & 0/27 \\
\bottomrule
\end{tabular}
\end{center}

\subsection{Family A --- isolated}

\paragraph{A:isolated} 11 eigenvalues in $[5, 300]$; sweep took 320\,s.

\begin{center}
\small
\begin{tabular}{lrrrrrrrr}
\toprule
& \multicolumn{2}{c}{$\pr=10^{-4}$} & \multicolumn{2}{c}{$10^{-8}$} & \multicolumn{2}{c}{$10^{-11}$} & \multicolumn{2}{c}{$10^{-14}$} \\
\cmidrule(lr){2-3}\cmidrule(lr){4-5}\cmidrule(lr){6-7}\cmidrule(lr){8-9}
method & orc & aut & orc & aut & orc & aut & orc & aut \\
\midrule
\texttt{uniform\_grid} & 830 & \textbf{--} & 830 & \textbf{--} & 830 & \textbf{--} & 830 & \textbf{--} \\
\texttt{lappy} & \textbf{--} & \textbf{--} & \textbf{--} & \textbf{--} & \textbf{--} & \textbf{--} & \textbf{--} & \textbf{--} \\
\texttt{lappy\_repaired} & 1449 & 1449 & 1449 & \textbf{--} & 1449 & \textbf{--} & 1449 & \textbf{--} \\
\texttt{lipschitz} & 610 & \textbf{--} & 610 & \textbf{--} & 610 & \textbf{--} & 610 & \textbf{--} \\
\texttt{swing\_cells} & \textbf{--} & \textbf{--} & \textbf{--} & \textbf{--} & \textbf{--} & \textbf{--} & \textbf{--} & \textbf{--} \\
\texttt{swing\_direct} & 923 & 923 & 1302 & 1302 & \textbf{--} & \textbf{--} & \textbf{--} & \textbf{--} \\
\midrule
\texttt{single\_bracket} & \textbf{--} & \textbf{--} & \textbf{--} & \textbf{--} & \textbf{--} & \textbf{--} & \textbf{--} & \textbf{--} \\
\texttt{fine\_grid} & \textbf{--} & \textbf{--} & \textbf{--} & \textbf{--} & \textbf{--} & \textbf{--} & \textbf{--} & \textbf{--} \\
\bottomrule
\end{tabular}
\end{center}

\subsection{Family B --- pair ladder}

\paragraph{B:d=1.8e-01} 4 eigenvalues in $[10, 88]$; sweep took 151\,s.

\begin{center}
\small
\begin{tabular}{lrrrrrrrr}
\toprule
& \multicolumn{2}{c}{$\pr=10^{-4}$} & \multicolumn{2}{c}{$10^{-8}$} & \multicolumn{2}{c}{$10^{-11}$} & \multicolumn{2}{c}{$10^{-14}$} \\
\cmidrule(lr){2-3}\cmidrule(lr){4-5}\cmidrule(lr){6-7}\cmidrule(lr){8-9}
method & orc & aut & orc & aut & orc & aut & orc & aut \\
\midrule
\texttt{uniform\_grid} & \textbf{--} & \textbf{--} & \textbf{--} & \textbf{--} & \textbf{--} & \textbf{--} & \textbf{--} & \textbf{--} \\
\texttt{lappy} & 1016 & \textbf{--} & 1016 & \textbf{--} & 1016 & \textbf{--} & 1016 & \textbf{--} \\
\texttt{lappy\_repaired} & 1103 & \textbf{--} & 1103 & \textbf{--} & 1103 & \textbf{--} & 1103 & \textbf{--} \\
\texttt{lipschitz} & 1298 & 1298 & 1298 & 1298 & 1298 & 1298 & 1298 & 1298 \\
\texttt{swing\_cells} & 528 & 528 & 788 & 788 & 988 & 988 & 1188 & 1188 \\
\texttt{swing\_direct} & 688 & 688 & 942 & 942 & 1150 & 1150 & 1352 & 1352 \\
\midrule
\texttt{single\_bracket} & \textbf{--} & \textbf{--} & \textbf{--} & \textbf{--} & \textbf{--} & \textbf{--} & \textbf{--} & \textbf{--} \\
\texttt{fine\_grid} & \textbf{--} & \textbf{--} & \textbf{--} & \textbf{--} & \textbf{--} & \textbf{--} & \textbf{--} & \textbf{--} \\
\bottomrule
\end{tabular}
\end{center}

\paragraph{B:d=1.8e-02} 4 eigenvalues in $[10, 88]$; sweep took 144\,s.

\begin{center}
\small
\begin{tabular}{lrrrrrrrr}
\toprule
& \multicolumn{2}{c}{$\pr=10^{-4}$} & \multicolumn{2}{c}{$10^{-8}$} & \multicolumn{2}{c}{$10^{-11}$} & \multicolumn{2}{c}{$10^{-14}$} \\
\cmidrule(lr){2-3}\cmidrule(lr){4-5}\cmidrule(lr){6-7}\cmidrule(lr){8-9}
method & orc & aut & orc & aut & orc & aut & orc & aut \\
\midrule
\texttt{uniform\_grid} & \textbf{--} & \textbf{--} & \textbf{--} & \textbf{--} & \textbf{--} & \textbf{--} & \textbf{--} & \textbf{--} \\
\texttt{lappy} & \textbf{--} & \textbf{--} & \textbf{--} & \textbf{--} & \textbf{--} & \textbf{--} & \textbf{--} & \textbf{--} \\
\texttt{lappy\_repaired} & 830 & \textbf{--} & 830 & \textbf{--} & 830 & \textbf{--} & 830 & \textbf{--} \\
\texttt{lipschitz} & \textbf{--} & \textbf{--} & \textbf{--} & \textbf{--} & \textbf{--} & \textbf{--} & \textbf{--} & \textbf{--} \\
\texttt{swing\_cells} & \textbf{--} & \textbf{--} & \textbf{--} & \textbf{--} & \textbf{--} & \textbf{--} & \textbf{--} & \textbf{--} \\
\texttt{swing\_direct} & \textbf{--} & \textbf{--} & \textbf{--} & \textbf{--} & \textbf{--} & \textbf{--} & \textbf{--} & \textbf{--} \\
\midrule
\texttt{single\_bracket} & \textbf{--} & \textbf{--} & \textbf{--} & \textbf{--} & \textbf{--} & \textbf{--} & \textbf{--} & \textbf{--} \\
\texttt{fine\_grid} & \textbf{--} & \textbf{--} & \textbf{--} & \textbf{--} & \textbf{--} & \textbf{--} & \textbf{--} & \textbf{--} \\
\bottomrule
\end{tabular}
\end{center}

\paragraph{B:d=1.8e-03} 4 eigenvalues in $[10, 88]$; sweep took 152\,s.

\begin{center}
\small
\begin{tabular}{lrrrrrrrr}
\toprule
& \multicolumn{2}{c}{$\pr=10^{-4}$} & \multicolumn{2}{c}{$10^{-8}$} & \multicolumn{2}{c}{$10^{-11}$} & \multicolumn{2}{c}{$10^{-14}$} \\
\cmidrule(lr){2-3}\cmidrule(lr){4-5}\cmidrule(lr){6-7}\cmidrule(lr){8-9}
method & orc & aut & orc & aut & orc & aut & orc & aut \\
\midrule
\texttt{uniform\_grid} & 86 & 158 & \textbf{--} & \textbf{--} & \textbf{--} & \textbf{--} & \textbf{--} & \textbf{--} \\
\texttt{lappy} & 122 & 302 & \textbf{--} & \textbf{--} & \textbf{--} & \textbf{--} & \textbf{--} & \textbf{--} \\
\texttt{lappy\_repaired} & \textbf{--} & \textbf{--} & \textbf{--} & \textbf{--} & \textbf{--} & \textbf{--} & \textbf{--} & \textbf{--} \\
\texttt{lipschitz} & 286 & 709 & \textbf{--} & \textbf{--} & \textbf{--} & \textbf{--} & \textbf{--} & \textbf{--} \\
\texttt{swing\_cells} & 184 & 704 & \textbf{--} & \textbf{--} & \textbf{--} & \textbf{--} & \textbf{--} & \textbf{--} \\
\texttt{swing\_direct} & 219 & 850 & \textbf{--} & \textbf{--} & \textbf{--} & \textbf{--} & \textbf{--} & \textbf{--} \\
\midrule
\texttt{single\_bracket} & \textbf{--} & \textbf{--} & \textbf{--} & \textbf{--} & \textbf{--} & \textbf{--} & \textbf{--} & \textbf{--} \\
\texttt{fine\_grid} & \textbf{--} & \textbf{--} & \textbf{--} & \textbf{--} & \textbf{--} & \textbf{--} & \textbf{--} & \textbf{--} \\
\bottomrule
\end{tabular}
\end{center}

\paragraph{B:d=1.8e-04} 4 eigenvalues in $[10, 88]$; sweep took 151\,s.

\begin{center}
\small
\begin{tabular}{lrrrrrrrr}
\toprule
& \multicolumn{2}{c}{$\pr=10^{-4}$} & \multicolumn{2}{c}{$10^{-8}$} & \multicolumn{2}{c}{$10^{-11}$} & \multicolumn{2}{c}{$10^{-14}$} \\
\cmidrule(lr){2-3}\cmidrule(lr){4-5}\cmidrule(lr){6-7}\cmidrule(lr){8-9}
method & orc & aut & orc & aut & orc & aut & orc & aut \\
\midrule
\texttt{uniform\_grid} & 86 & 158 & \textbf{--} & \textbf{--} & \textbf{--} & \textbf{--} & \textbf{--} & \textbf{--} \\
\texttt{lappy} & 122 & 302 & \textbf{--} & \textbf{--} & \textbf{--} & \textbf{--} & \textbf{--} & \textbf{--} \\
\texttt{lappy\_repaired} & 1454 & \textbf{--} & \textbf{--} & \textbf{--} & \textbf{--} & \textbf{--} & \textbf{--} & \textbf{--} \\
\texttt{lipschitz} & 289 & 1174 & \textbf{--} & \textbf{--} & \textbf{--} & \textbf{--} & \textbf{--} & \textbf{--} \\
\texttt{swing\_cells} & 184 & 648 & \textbf{--} & \textbf{--} & \textbf{--} & \textbf{--} & \textbf{--} & \textbf{--} \\
\texttt{swing\_direct} & 227 & 779 & \textbf{--} & \textbf{--} & \textbf{--} & \textbf{--} & \textbf{--} & \textbf{--} \\
\midrule
\texttt{single\_bracket} & \textbf{--} & \textbf{--} & \textbf{--} & \textbf{--} & \textbf{--} & \textbf{--} & \textbf{--} & \textbf{--} \\
\texttt{fine\_grid} & \textbf{--} & \textbf{--} & \textbf{--} & \textbf{--} & \textbf{--} & \textbf{--} & \textbf{--} & \textbf{--} \\
\bottomrule
\end{tabular}
\end{center}

\paragraph{B:d=1.8e-08} 4 eigenvalues in $[10, 88]$; sweep took 166\,s.

\begin{center}
\small
\begin{tabular}{lrrrrrrrr}
\toprule
& \multicolumn{2}{c}{$\pr=10^{-4}$} & \multicolumn{2}{c}{$10^{-8}$} & \multicolumn{2}{c}{$10^{-11}$} & \multicolumn{2}{c}{$10^{-14}$} \\
\cmidrule(lr){2-3}\cmidrule(lr){4-5}\cmidrule(lr){6-7}\cmidrule(lr){8-9}
method & orc & aut & orc & aut & orc & aut & orc & aut \\
\midrule
\texttt{uniform\_grid} & 86 & 158 & 86 & 158 & \textbf{--} & \textbf{--} & \textbf{--} & \textbf{--} \\
\texttt{lappy} & 122 & 302 & 122 & 302 & \textbf{--} & \textbf{--} & \textbf{--} & \textbf{--} \\
\texttt{lappy\_repaired} & 1011 & \textbf{--} & 1011 & \textbf{--} & \textbf{--} & \textbf{--} & \textbf{--} & \textbf{--} \\
\texttt{lipschitz} & 441 & 1191 & 441 & 1191 & \textbf{--} & \textbf{--} & \textbf{--} & \textbf{--} \\
\texttt{swing\_cells} & 163 & 686 & 215 & 1063 & \textbf{--} & \textbf{--} & \textbf{--} & \textbf{--} \\
\texttt{swing\_direct} & 197 & 891 & 245 & 1257 & \textbf{--} & \textbf{--} & \textbf{--} & \textbf{--} \\
\bottomrule
\end{tabular}
\end{center}

\paragraph{B:d=1.8e-12} 4 eigenvalues in $[10, 88]$; sweep took 180\,s.

\begin{center}
\small
\begin{tabular}{lrrrrrrrr}
\toprule
& \multicolumn{2}{c}{$\pr=10^{-4}$} & \multicolumn{2}{c}{$10^{-8}$} & \multicolumn{2}{c}{$10^{-11}$} & \multicolumn{2}{c}{$10^{-14}$} \\
\cmidrule(lr){2-3}\cmidrule(lr){4-5}\cmidrule(lr){6-7}\cmidrule(lr){8-9}
method & orc & aut & orc & aut & orc & aut & orc & aut \\
\midrule
\texttt{uniform\_grid} & 86 & 158 & 86 & 158 & 86 & 158 & \textbf{--} & \textbf{--} \\
\texttt{lappy} & 122 & 302 & 122 & 302 & 122 & 302 & \textbf{--} & \textbf{--} \\
\texttt{lappy\_repaired} & \textbf{--} & \textbf{--} & \textbf{--} & \textbf{--} & \textbf{--} & \textbf{--} & \textbf{--} & \textbf{--} \\
\texttt{lipschitz} & 471 & 1246 & 471 & 1246 & 471 & 1246 & \textbf{--} & \textbf{--} \\
\texttt{swing\_cells} & 131 & 349 & 170 & 492 & 200 & 602 & \textbf{--} & \textbf{--} \\
\texttt{swing\_direct} & 152 & 418 & 191 & 561 & 215 & 661 & \textbf{--} & \textbf{--} \\
\bottomrule
\end{tabular}
\end{center}

\paragraph{B:d=5.8e-01} 4 eigenvalues in $[10, 88]$; sweep took 150\,s.

\begin{center}
\small
\begin{tabular}{lrrrrrrrr}
\toprule
& \multicolumn{2}{c}{$\pr=10^{-4}$} & \multicolumn{2}{c}{$10^{-8}$} & \multicolumn{2}{c}{$10^{-11}$} & \multicolumn{2}{c}{$10^{-14}$} \\
\cmidrule(lr){2-3}\cmidrule(lr){4-5}\cmidrule(lr){6-7}\cmidrule(lr){8-9}
method & orc & aut & orc & aut & orc & aut & orc & aut \\
\midrule
\texttt{uniform\_grid} & 446 & \textbf{--} & 446 & \textbf{--} & 446 & \textbf{--} & 446 & \textbf{--} \\
\texttt{lappy} & 302 & 302 & 302 & 302 & 302 & 302 & 302 & 302 \\
\texttt{lappy\_repaired} & 1316 & \textbf{--} & 1316 & \textbf{--} & 1316 & \textbf{--} & 1316 & \textbf{--} \\
\texttt{lipschitz} & 293 & 1193 & 293 & 1193 & 293 & 1193 & 293 & 1193 \\
\texttt{swing\_cells} & 131 & 379 & 170 & 548 & 200 & 678 & 230 & 808 \\
\texttt{swing\_direct} & 152 & 458 & 197 & 633 & 225 & 773 & 251 & 900 \\
\midrule
\texttt{single\_bracket} & \textbf{--} & \textbf{--} & \textbf{--} & \textbf{--} & \textbf{--} & \textbf{--} & \textbf{--} & \textbf{--} \\
\texttt{fine\_grid} & \textbf{--} & \textbf{--} & \textbf{--} & \textbf{--} & \textbf{--} & \textbf{--} & \textbf{--} & \textbf{--} \\
\bottomrule
\end{tabular}
\end{center}

\paragraph{B:d=5.9e-02} 4 eigenvalues in $[10, 88]$; sweep took 142\,s.

\begin{center}
\small
\begin{tabular}{lrrrrrrrr}
\toprule
& \multicolumn{2}{c}{$\pr=10^{-4}$} & \multicolumn{2}{c}{$10^{-8}$} & \multicolumn{2}{c}{$10^{-11}$} & \multicolumn{2}{c}{$10^{-14}$} \\
\cmidrule(lr){2-3}\cmidrule(lr){4-5}\cmidrule(lr){6-7}\cmidrule(lr){8-9}
method & orc & aut & orc & aut & orc & aut & orc & aut \\
\midrule
\texttt{uniform\_grid} & \textbf{--} & \textbf{--} & \textbf{--} & \textbf{--} & \textbf{--} & \textbf{--} & \textbf{--} & \textbf{--} \\
\texttt{lappy} & \textbf{--} & \textbf{--} & \textbf{--} & \textbf{--} & \textbf{--} & \textbf{--} & \textbf{--} & \textbf{--} \\
\texttt{lappy\_repaired} & 922 & \textbf{--} & 922 & \textbf{--} & 922 & \textbf{--} & 922 & \textbf{--} \\
\texttt{lipschitz} & 1308 & 1308 & 1308 & 1308 & 1308 & 1308 & 1308 & 1308 \\
\texttt{swing\_cells} & 184 & \textbf{--} & 249 & \textbf{--} & 299 & \textbf{--} & 349 & \textbf{--} \\
\texttt{swing\_direct} & 225 & \textbf{--} & 300 & \textbf{--} & 338 & \textbf{--} & 392 & \textbf{--} \\
\midrule
\texttt{single\_bracket} & \textbf{--} & \textbf{--} & \textbf{--} & \textbf{--} & \textbf{--} & \textbf{--} & \textbf{--} & \textbf{--} \\
\texttt{fine\_grid} & \textbf{--} & \textbf{--} & \textbf{--} & \textbf{--} & \textbf{--} & \textbf{--} & \textbf{--} & \textbf{--} \\
\bottomrule
\end{tabular}
\end{center}

\paragraph{B:d=5.9e-03} 4 eigenvalues in $[10, 88]$; sweep took 152\,s.

\begin{center}
\small
\begin{tabular}{lrrrrrrrr}
\toprule
& \multicolumn{2}{c}{$\pr=10^{-4}$} & \multicolumn{2}{c}{$10^{-8}$} & \multicolumn{2}{c}{$10^{-11}$} & \multicolumn{2}{c}{$10^{-14}$} \\
\cmidrule(lr){2-3}\cmidrule(lr){4-5}\cmidrule(lr){6-7}\cmidrule(lr){8-9}
method & orc & aut & orc & aut & orc & aut & orc & aut \\
\midrule
\texttt{uniform\_grid} & \textbf{--} & \textbf{--} & \textbf{--} & \textbf{--} & \textbf{--} & \textbf{--} & \textbf{--} & \textbf{--} \\
\texttt{lappy} & \textbf{--} & \textbf{--} & \textbf{--} & \textbf{--} & \textbf{--} & \textbf{--} & \textbf{--} & \textbf{--} \\
\texttt{lappy\_repaired} & 1426 & 1426 & \textbf{--} & \textbf{--} & \textbf{--} & \textbf{--} & \textbf{--} & \textbf{--} \\
\texttt{lipschitz} & \textbf{--} & \textbf{--} & \textbf{--} & \textbf{--} & \textbf{--} & \textbf{--} & \textbf{--} & \textbf{--} \\
\texttt{swing\_cells} & 658 & 658 & 1009 & 1009 & 1279 & 1279 & \textbf{--} & \textbf{--} \\
\texttt{swing\_direct} & 844 & 844 & 1213 & 1213 & 1462 & 1462 & \textbf{--} & \textbf{--} \\
\midrule
\texttt{single\_bracket} & \textbf{--} & \textbf{--} & \textbf{--} & \textbf{--} & \textbf{--} & \textbf{--} & \textbf{--} & \textbf{--} \\
\texttt{fine\_grid} & \textbf{--} & \textbf{--} & \textbf{--} & \textbf{--} & \textbf{--} & \textbf{--} & \textbf{--} & \textbf{--} \\
\bottomrule
\end{tabular}
\end{center}

\paragraph{B:d=5.9e-04} 4 eigenvalues in $[10, 88]$; sweep took 149\,s.

\begin{center}
\small
\begin{tabular}{lrrrrrrrr}
\toprule
& \multicolumn{2}{c}{$\pr=10^{-4}$} & \multicolumn{2}{c}{$10^{-8}$} & \multicolumn{2}{c}{$10^{-11}$} & \multicolumn{2}{c}{$10^{-14}$} \\
\cmidrule(lr){2-3}\cmidrule(lr){4-5}\cmidrule(lr){6-7}\cmidrule(lr){8-9}
method & orc & aut & orc & aut & orc & aut & orc & aut \\
\midrule
\texttt{uniform\_grid} & 86 & 158 & \textbf{--} & \textbf{--} & \textbf{--} & \textbf{--} & \textbf{--} & \textbf{--} \\
\texttt{lappy} & 122 & 302 & \textbf{--} & \textbf{--} & \textbf{--} & \textbf{--} & \textbf{--} & \textbf{--} \\
\texttt{lappy\_repaired} & 684 & 1324 & \textbf{--} & \textbf{--} & \textbf{--} & \textbf{--} & \textbf{--} & \textbf{--} \\
\texttt{lipschitz} & 305 & 1283 & \textbf{--} & \textbf{--} & \textbf{--} & \textbf{--} & \textbf{--} & \textbf{--} \\
\texttt{swing\_cells} & 137 & 510 & \textbf{--} & \textbf{--} & \textbf{--} & \textbf{--} & \textbf{--} & \textbf{--} \\
\texttt{swing\_direct} & 159 & 648 & \textbf{--} & \textbf{--} & \textbf{--} & \textbf{--} & \textbf{--} & \textbf{--} \\
\midrule
\texttt{single\_bracket} & \textbf{--} & \textbf{--} & \textbf{--} & \textbf{--} & \textbf{--} & \textbf{--} & \textbf{--} & \textbf{--} \\
\texttt{fine\_grid} & \textbf{--} & \textbf{--} & \textbf{--} & \textbf{--} & \textbf{--} & \textbf{--} & \textbf{--} & \textbf{--} \\
\bottomrule
\end{tabular}
\end{center}

\paragraph{B:d=5.9e-05} 4 eigenvalues in $[10, 88]$; sweep took 150\,s.

\begin{center}
\small
\begin{tabular}{lrrrrrrrr}
\toprule
& \multicolumn{2}{c}{$\pr=10^{-4}$} & \multicolumn{2}{c}{$10^{-8}$} & \multicolumn{2}{c}{$10^{-11}$} & \multicolumn{2}{c}{$10^{-14}$} \\
\cmidrule(lr){2-3}\cmidrule(lr){4-5}\cmidrule(lr){6-7}\cmidrule(lr){8-9}
method & orc & aut & orc & aut & orc & aut & orc & aut \\
\midrule
\texttt{uniform\_grid} & 86 & 158 & \textbf{--} & \textbf{--} & \textbf{--} & \textbf{--} & \textbf{--} & \textbf{--} \\
\texttt{lappy} & 122 & 302 & \textbf{--} & \textbf{--} & \textbf{--} & \textbf{--} & \textbf{--} & \textbf{--} \\
\texttt{lappy\_repaired} & 1223 & \textbf{--} & \textbf{--} & \textbf{--} & \textbf{--} & \textbf{--} & \textbf{--} & \textbf{--} \\
\texttt{lipschitz} & 503 & 1349 & \textbf{--} & \textbf{--} & \textbf{--} & \textbf{--} & \textbf{--} & \textbf{--} \\
\texttt{swing\_cells} & 172 & 556 & \textbf{--} & \textbf{--} & \textbf{--} & \textbf{--} & \textbf{--} & \textbf{--} \\
\texttt{swing\_direct} & 216 & 694 & \textbf{--} & \textbf{--} & \textbf{--} & \textbf{--} & \textbf{--} & \textbf{--} \\
\midrule
\texttt{single\_bracket} & \textbf{--} & \textbf{--} & \textbf{--} & \textbf{--} & \textbf{--} & \textbf{--} & \textbf{--} & \textbf{--} \\
\texttt{fine\_grid} & \textbf{--} & \textbf{--} & \textbf{--} & \textbf{--} & \textbf{--} & \textbf{--} & \textbf{--} & \textbf{--} \\
\bottomrule
\end{tabular}
\end{center}

\paragraph{B:d=5.9e-07} 4 eigenvalues in $[10, 88]$; sweep took 173\,s.

\begin{center}
\small
\begin{tabular}{lrrrrrrrr}
\toprule
& \multicolumn{2}{c}{$\pr=10^{-4}$} & \multicolumn{2}{c}{$10^{-8}$} & \multicolumn{2}{c}{$10^{-11}$} & \multicolumn{2}{c}{$10^{-14}$} \\
\cmidrule(lr){2-3}\cmidrule(lr){4-5}\cmidrule(lr){6-7}\cmidrule(lr){8-9}
method & orc & aut & orc & aut & orc & aut & orc & aut \\
\midrule
\texttt{uniform\_grid} & 86 & 158 & \textbf{--} & \textbf{--} & \textbf{--} & \textbf{--} & \textbf{--} & \textbf{--} \\
\texttt{lappy} & 122 & 302 & \textbf{--} & \textbf{--} & \textbf{--} & \textbf{--} & \textbf{--} & \textbf{--} \\
\texttt{lappy\_repaired} & 1025 & \textbf{--} & \textbf{--} & \textbf{--} & \textbf{--} & \textbf{--} & \textbf{--} & \textbf{--} \\
\texttt{lipschitz} & 502 & 1379 & \textbf{--} & \textbf{--} & \textbf{--} & \textbf{--} & \textbf{--} & \textbf{--} \\
\texttt{swing\_cells} & 171 & 482 & \textbf{--} & \textbf{--} & \textbf{--} & \textbf{--} & \textbf{--} & \textbf{--} \\
\texttt{swing\_direct} & 208 & 617 & \textbf{--} & \textbf{--} & \textbf{--} & \textbf{--} & \textbf{--} & \textbf{--} \\
\bottomrule
\end{tabular}
\end{center}

\paragraph{B:d=5.9e-10} 4 eigenvalues in $[10, 88]$; sweep took 166\,s.

\begin{center}
\small
\begin{tabular}{lrrrrrrrr}
\toprule
& \multicolumn{2}{c}{$\pr=10^{-4}$} & \multicolumn{2}{c}{$10^{-8}$} & \multicolumn{2}{c}{$10^{-11}$} & \multicolumn{2}{c}{$10^{-14}$} \\
\cmidrule(lr){2-3}\cmidrule(lr){4-5}\cmidrule(lr){6-7}\cmidrule(lr){8-9}
method & orc & aut & orc & aut & orc & aut & orc & aut \\
\midrule
\texttt{uniform\_grid} & 86 & 158 & 86 & 158 & \textbf{--} & \textbf{--} & \textbf{--} & \textbf{--} \\
\texttt{lappy} & 122 & 302 & 122 & 302 & \textbf{--} & \textbf{--} & \textbf{--} & \textbf{--} \\
\texttt{lappy\_repaired} & 1295 & \textbf{--} & 1295 & \textbf{--} & \textbf{--} & \textbf{--} & \textbf{--} & \textbf{--} \\
\texttt{lipschitz} & 306 & 789 & 306 & 789 & \textbf{--} & \textbf{--} & \textbf{--} & \textbf{--} \\
\texttt{swing\_cells} & 195 & 644 & 273 & 982 & \textbf{--} & \textbf{--} & \textbf{--} & \textbf{--} \\
\texttt{swing\_direct} & 235 & 804 & 319 & 1149 & \textbf{--} & \textbf{--} & \textbf{--} & \textbf{--} \\
\bottomrule
\end{tabular}
\end{center}

\paragraph{B:d=5.9e-12} 4 eigenvalues in $[10, 88]$; sweep took 189\,s.

\begin{center}
\small
\begin{tabular}{lrrrrrrrr}
\toprule
& \multicolumn{2}{c}{$\pr=10^{-4}$} & \multicolumn{2}{c}{$10^{-8}$} & \multicolumn{2}{c}{$10^{-11}$} & \multicolumn{2}{c}{$10^{-14}$} \\
\cmidrule(lr){2-3}\cmidrule(lr){4-5}\cmidrule(lr){6-7}\cmidrule(lr){8-9}
method & orc & aut & orc & aut & orc & aut & orc & aut \\
\midrule
\texttt{uniform\_grid} & 86 & 158 & 86 & 158 & 86 & 158 & \textbf{--} & \textbf{--} \\
\texttt{lappy} & 122 & 302 & 122 & 302 & 122 & 302 & \textbf{--} & \textbf{--} \\
\texttt{lappy\_repaired} & 583 & 1063 & 583 & \textbf{--} & 583 & \textbf{--} & \textbf{--} & \textbf{--} \\
\texttt{lipschitz} & 488 & 1347 & 488 & 1347 & 488 & 1347 & \textbf{--} & \textbf{--} \\
\texttt{swing\_cells} & 179 & 686 & 244 & 1063 & 294 & 1353 & \textbf{--} & \textbf{--} \\
\texttt{swing\_direct} & 218 & 887 & 283 & 1257 & 333 & \textbf{--} & \textbf{--} & \textbf{--} \\
\bottomrule
\end{tabular}
\end{center}

\subsection{Family C --- clusters of 3 and 4}

\paragraph{C:m=3,s=9.3e+00} 3 eigenvalues in $[478.7, 508.3]$; sweep took 800\,s.

\begin{center}
\small
\begin{tabular}{lrrrrrrrr}
\toprule
& \multicolumn{2}{c}{$\pr=10^{-4}$} & \multicolumn{2}{c}{$10^{-8}$} & \multicolumn{2}{c}{$10^{-11}$} & \multicolumn{2}{c}{$10^{-14}$} \\
\cmidrule(lr){2-3}\cmidrule(lr){4-5}\cmidrule(lr){6-7}\cmidrule(lr){8-9}
method & orc & aut & orc & aut & orc & aut & orc & aut \\
\midrule
\texttt{uniform\_grid} & 94 & 190 & 94 & 190 & 94 & 190 & 94 & 190 \\
\texttt{lappy} & 142 & 382 & 142 & 382 & 142 & 382 & 142 & 382 \\
\texttt{lappy\_repaired} & 278 & 278 & 634 & \textbf{--} & 634 & \textbf{--} & 634 & \textbf{--} \\
\texttt{lipschitz} & 263 & 403 & 263 & 663 & 263 & 663 & 263 & 663 \\
\texttt{swing\_cells} & 119 & 327 & 158 & 496 & 188 & 626 & 218 & 756 \\
\texttt{swing\_direct} & 142 & 400 & 189 & 567 & 213 & 705 & 245 & 837 \\
\midrule
\texttt{single\_bracket} & \textbf{--} & \textbf{--} & \textbf{--} & \textbf{--} & \textbf{--} & \textbf{--} & \textbf{--} & \textbf{--} \\
\texttt{fine\_grid} & \textbf{--} & \textbf{--} & \textbf{--} & \textbf{--} & \textbf{--} & \textbf{--} & \textbf{--} & \textbf{--} \\
\bottomrule
\end{tabular}
\end{center}

\paragraph{C:m=3,s=9.5e-01} 3 eigenvalues in $[478.7, 508.3]$; sweep took 670\,s.

\begin{center}
\small
\begin{tabular}{lrrrrrrrr}
\toprule
& \multicolumn{2}{c}{$\pr=10^{-4}$} & \multicolumn{2}{c}{$10^{-8}$} & \multicolumn{2}{c}{$10^{-11}$} & \multicolumn{2}{c}{$10^{-14}$} \\
\cmidrule(lr){2-3}\cmidrule(lr){4-5}\cmidrule(lr){6-7}\cmidrule(lr){8-9}
method & orc & aut & orc & aut & orc & aut & orc & aut \\
\midrule
\texttt{uniform\_grid} & 190 & \textbf{--} & 190 & \textbf{--} & 190 & \textbf{--} & 190 & \textbf{--} \\
\texttt{lappy} & 382 & \textbf{--} & 382 & \textbf{--} & 382 & \textbf{--} & 382 & \textbf{--} \\
\texttt{lappy\_repaired} & 252 & 252 & 387 & 1352 & 387 & 1352 & 387 & 1352 \\
\texttt{lipschitz} & 162 & 347 & 162 & 347 & 162 & 347 & 162 & 347 \\
\texttt{swing\_cells} & 1312 & \textbf{--} & \textbf{--} & \textbf{--} & \textbf{--} & \textbf{--} & \textbf{--} & \textbf{--} \\
\texttt{swing\_direct} & \textbf{--} & \textbf{--} & \textbf{--} & \textbf{--} & \textbf{--} & \textbf{--} & \textbf{--} & \textbf{--} \\
\midrule
\texttt{single\_bracket} & \textbf{--} & \textbf{--} & \textbf{--} & \textbf{--} & \textbf{--} & \textbf{--} & \textbf{--} & \textbf{--} \\
\texttt{fine\_grid} & \textbf{--} & \textbf{--} & \textbf{--} & \textbf{--} & \textbf{--} & \textbf{--} & \textbf{--} & \textbf{--} \\
\bottomrule
\end{tabular}
\end{center}

\paragraph{C:m=3,s=9.5e-02} 3 eigenvalues in $[478.7, 508.3]$; sweep took 503\,s.

\begin{center}
\small
\begin{tabular}{lrrrrrrrr}
\toprule
& \multicolumn{2}{c}{$\pr=10^{-4}$} & \multicolumn{2}{c}{$10^{-8}$} & \multicolumn{2}{c}{$10^{-11}$} & \multicolumn{2}{c}{$10^{-14}$} \\
\cmidrule(lr){2-3}\cmidrule(lr){4-5}\cmidrule(lr){6-7}\cmidrule(lr){8-9}
method & orc & aut & orc & aut & orc & aut & orc & aut \\
\midrule
\texttt{uniform\_grid} & \textbf{--} & \textbf{--} & \textbf{--} & \textbf{--} & \textbf{--} & \textbf{--} & \textbf{--} & \textbf{--} \\
\texttt{lappy} & \textbf{--} & \textbf{--} & \textbf{--} & \textbf{--} & \textbf{--} & \textbf{--} & \textbf{--} & \textbf{--} \\
\texttt{lappy\_repaired} & 165 & \textbf{--} & \textbf{--} & \textbf{--} & \textbf{--} & \textbf{--} & \textbf{--} & \textbf{--} \\
\texttt{lipschitz} & \textbf{--} & \textbf{--} & \textbf{--} & \textbf{--} & \textbf{--} & \textbf{--} & \textbf{--} & \textbf{--} \\
\texttt{swing\_cells} & 102 & 262 & \textbf{--} & \textbf{--} & \textbf{--} & \textbf{--} & \textbf{--} & \textbf{--} \\
\texttt{swing\_direct} & 122 & 336 & \textbf{--} & \textbf{--} & \textbf{--} & \textbf{--} & \textbf{--} & \textbf{--} \\
\midrule
\texttt{single\_bracket} & \textbf{--} & \textbf{--} & \textbf{--} & \textbf{--} & \textbf{--} & \textbf{--} & \textbf{--} & \textbf{--} \\
\texttt{fine\_grid} & \textbf{--} & \textbf{--} & \textbf{--} & \textbf{--} & \textbf{--} & \textbf{--} & \textbf{--} & \textbf{--} \\
\bottomrule
\end{tabular}
\end{center}

\paragraph{C:m=3,s=9.5e-03} 3 eigenvalues in $[478.7, 508.3]$; sweep took 576\,s.

\begin{center}
\small
\begin{tabular}{lrrrrrrrr}
\toprule
& \multicolumn{2}{c}{$\pr=10^{-4}$} & \multicolumn{2}{c}{$10^{-8}$} & \multicolumn{2}{c}{$10^{-11}$} & \multicolumn{2}{c}{$10^{-14}$} \\
\cmidrule(lr){2-3}\cmidrule(lr){4-5}\cmidrule(lr){6-7}\cmidrule(lr){8-9}
method & orc & aut & orc & aut & orc & aut & orc & aut \\
\midrule
\texttt{uniform\_grid} & 76 & 94 & \textbf{--} & \textbf{--} & \textbf{--} & \textbf{--} & \textbf{--} & \textbf{--} \\
\texttt{lappy} & 94 & 142 & \textbf{--} & \textbf{--} & \textbf{--} & \textbf{--} & \textbf{--} & \textbf{--} \\
\texttt{lappy\_repaired} & \textbf{--} & \textbf{--} & 899 & 899 & 899 & 899 & 899 & 899 \\
\texttt{lipschitz} & 160 & 344 & \textbf{--} & \textbf{--} & \textbf{--} & \textbf{--} & \textbf{--} & \textbf{--} \\
\texttt{swing\_cells} & 138 & 346 & \textbf{--} & \textbf{--} & \textbf{--} & \textbf{--} & \textbf{--} & \textbf{--} \\
\texttt{swing\_direct} & 162 & 429 & \textbf{--} & \textbf{--} & \textbf{--} & \textbf{--} & \textbf{--} & \textbf{--} \\
\midrule
\texttt{single\_bracket} & 1 & 1 & \textbf{--} & \textbf{--} & \textbf{--} & \textbf{--} & \textbf{--} & \textbf{--} \\
\texttt{fine\_grid} & \textbf{--} & \textbf{--} & \textbf{--} & \textbf{--} & \textbf{--} & \textbf{--} & \textbf{--} & \textbf{--} \\
\bottomrule
\end{tabular}
\end{center}

\paragraph{C:m=4,s=1.2e+00} 4 eigenvalues in $[622.3, 660.8]$; sweep took 792\,s.

\begin{center}
\small
\begin{tabular}{lrrrrrrrr}
\toprule
& \multicolumn{2}{c}{$\pr=10^{-4}$} & \multicolumn{2}{c}{$10^{-8}$} & \multicolumn{2}{c}{$10^{-11}$} & \multicolumn{2}{c}{$10^{-14}$} \\
\cmidrule(lr){2-3}\cmidrule(lr){4-5}\cmidrule(lr){6-7}\cmidrule(lr){8-9}
method & orc & aut & orc & aut & orc & aut & orc & aut \\
\midrule
\texttt{uniform\_grid} & \textbf{--} & \textbf{--} & \textbf{--} & \textbf{--} & \textbf{--} & \textbf{--} & \textbf{--} & \textbf{--} \\
\texttt{lappy} & 382 & \textbf{--} & 382 & \textbf{--} & 382 & \textbf{--} & 382 & \textbf{--} \\
\texttt{lappy\_repaired} & 328 & 328 & 625 & \textbf{--} & 625 & \textbf{--} & 625 & \textbf{--} \\
\texttt{lipschitz} & 170 & 385 & 170 & 385 & 170 & 385 & 170 & 385 \\
\texttt{swing\_cells} & 378 & 378 & 547 & 547 & 677 & 677 & 807 & 807 \\
\texttt{swing\_direct} & 491 & 491 & 668 & 668 & 788 & 788 & 928 & 928 \\
\midrule
\texttt{single\_bracket} & \textbf{--} & \textbf{--} & \textbf{--} & \textbf{--} & \textbf{--} & \textbf{--} & \textbf{--} & \textbf{--} \\
\texttt{fine\_grid} & \textbf{--} & \textbf{--} & \textbf{--} & \textbf{--} & \textbf{--} & \textbf{--} & \textbf{--} & \textbf{--} \\
\bottomrule
\end{tabular}
\end{center}

\paragraph{C:m=4,s=1.2e+01} 4 eigenvalues in $[622.3, 660.8]$; sweep took 1062\,s.

\begin{center}
\small
\begin{tabular}{lrrrrrrrr}
\toprule
& \multicolumn{2}{c}{$\pr=10^{-4}$} & \multicolumn{2}{c}{$10^{-8}$} & \multicolumn{2}{c}{$10^{-11}$} & \multicolumn{2}{c}{$10^{-14}$} \\
\cmidrule(lr){2-3}\cmidrule(lr){4-5}\cmidrule(lr){6-7}\cmidrule(lr){8-9}
method & orc & aut & orc & aut & orc & aut & orc & aut \\
\midrule
\texttt{uniform\_grid} & 94 & 190 & 94 & 190 & 94 & 190 & 94 & 190 \\
\texttt{lappy} & 142 & 382 & 142 & 382 & 142 & 382 & 142 & 382 \\
\texttt{lappy\_repaired} & 334 & 334 & 1054 & \textbf{--} & 1054 & \textbf{--} & 1054 & \textbf{--} \\
\texttt{lipschitz} & 273 & 442 & 273 & 762 & 273 & 762 & 273 & 762 \\
\texttt{swing\_cells} & 144 & 438 & 196 & 685 & 236 & 875 & 276 & 1065 \\
\texttt{swing\_direct} & 180 & 551 & 242 & 808 & 272 & 986 & 314 & 1188 \\
\midrule
\texttt{single\_bracket} & \textbf{--} & \textbf{--} & \textbf{--} & \textbf{--} & \textbf{--} & \textbf{--} & \textbf{--} & \textbf{--} \\
\texttt{fine\_grid} & \textbf{--} & \textbf{--} & \textbf{--} & \textbf{--} & \textbf{--} & \textbf{--} & \textbf{--} & \textbf{--} \\
\bottomrule
\end{tabular}
\end{center}

\paragraph{C:m=4,s=1.2e-01} 4 eigenvalues in $[622.3, 660.8]$; sweep took 794\,s.

\begin{center}
\small
\begin{tabular}{lrrrrrrrr}
\toprule
& \multicolumn{2}{c}{$\pr=10^{-4}$} & \multicolumn{2}{c}{$10^{-8}$} & \multicolumn{2}{c}{$10^{-11}$} & \multicolumn{2}{c}{$10^{-14}$} \\
\cmidrule(lr){2-3}\cmidrule(lr){4-5}\cmidrule(lr){6-7}\cmidrule(lr){8-9}
method & orc & aut & orc & aut & orc & aut & orc & aut \\
\midrule
\texttt{uniform\_grid} & \textbf{--} & \textbf{--} & \textbf{--} & \textbf{--} & \textbf{--} & \textbf{--} & \textbf{--} & \textbf{--} \\
\texttt{lappy} & \textbf{--} & \textbf{--} & \textbf{--} & \textbf{--} & \textbf{--} & \textbf{--} & \textbf{--} & \textbf{--} \\
\texttt{lappy\_repaired} & 478 & \textbf{--} & 752 & \textbf{--} & 752 & \textbf{--} & 752 & \textbf{--} \\
\texttt{lipschitz} & 407 & 407 & 772 & \textbf{--} & 772 & \textbf{--} & 772 & \textbf{--} \\
\texttt{swing\_cells} & 1145 & \textbf{--} & \textbf{--} & \textbf{--} & \textbf{--} & \textbf{--} & \textbf{--} & \textbf{--} \\
\texttt{swing\_direct} & \textbf{--} & \textbf{--} & \textbf{--} & \textbf{--} & \textbf{--} & \textbf{--} & \textbf{--} & \textbf{--} \\
\midrule
\texttt{single\_bracket} & \textbf{--} & \textbf{--} & \textbf{--} & \textbf{--} & \textbf{--} & \textbf{--} & \textbf{--} & \textbf{--} \\
\texttt{fine\_grid} & \textbf{--} & \textbf{--} & \textbf{--} & \textbf{--} & \textbf{--} & \textbf{--} & \textbf{--} & \textbf{--} \\
\bottomrule
\end{tabular}
\end{center}

\paragraph{C:m=4,s=1.2e-02} 4 eigenvalues in $[622.3, 660.8]$; sweep took 758\,s.

\begin{center}
\small
\begin{tabular}{lrrrrrrrr}
\toprule
& \multicolumn{2}{c}{$\pr=10^{-4}$} & \multicolumn{2}{c}{$10^{-8}$} & \multicolumn{2}{c}{$10^{-11}$} & \multicolumn{2}{c}{$10^{-14}$} \\
\cmidrule(lr){2-3}\cmidrule(lr){4-5}\cmidrule(lr){6-7}\cmidrule(lr){8-9}
method & orc & aut & orc & aut & orc & aut & orc & aut \\
\midrule
\texttt{uniform\_grid} & 76 & 94 & \textbf{--} & \textbf{--} & \textbf{--} & \textbf{--} & \textbf{--} & \textbf{--} \\
\texttt{lappy} & 94 & 142 & \textbf{--} & \textbf{--} & \textbf{--} & \textbf{--} & \textbf{--} & \textbf{--} \\
\texttt{lappy\_repaired} & 381 & 381 & \textbf{--} & \textbf{--} & \textbf{--} & \textbf{--} & \textbf{--} & \textbf{--} \\
\texttt{lipschitz} & 160 & 364 & \textbf{--} & \textbf{--} & \textbf{--} & \textbf{--} & \textbf{--} & \textbf{--} \\
\texttt{swing\_cells} & 131 & 338 & \textbf{--} & \textbf{--} & \textbf{--} & \textbf{--} & \textbf{--} & \textbf{--} \\
\texttt{swing\_direct} & 154 & 389 & \textbf{--} & \textbf{--} & \textbf{--} & \textbf{--} & \textbf{--} & \textbf{--} \\
\midrule
\texttt{single\_bracket} & 1 & 1 & \textbf{--} & \textbf{--} & \textbf{--} & \textbf{--} & \textbf{--} & \textbf{--} \\
\texttt{fine\_grid} & \textbf{--} & \textbf{--} & \textbf{--} & \textbf{--} & \textbf{--} & \textbf{--} & \textbf{--} & \textbf{--} \\
\bottomrule
\end{tabular}
\end{center}

\subsection{Family D --- exact degeneracy}

\paragraph{D:mult=2} 2 eigenvalues in $[239.3, 254.1]$; sweep took 345\,s.

\begin{center}
\small
\begin{tabular}{lrrrrrrrr}
\toprule
& \multicolumn{2}{c}{$\pr=10^{-4}$} & \multicolumn{2}{c}{$10^{-8}$} & \multicolumn{2}{c}{$10^{-11}$} & \multicolumn{2}{c}{$10^{-14}$} \\
\cmidrule(lr){2-3}\cmidrule(lr){4-5}\cmidrule(lr){6-7}\cmidrule(lr){8-9}
method & orc & aut & orc & aut & orc & aut & orc & aut \\
\midrule
\texttt{uniform\_grid} & 76 & 94 & 76 & 94 & 76 & 94 & 76 & 94 \\
\texttt{lappy} & 94 & 142 & 94 & 142 & 94 & 142 & 94 & 142 \\
\texttt{lappy\_repaired} & 210 & 210 & 295 & 936 & 295 & 936 & 295 & 936 \\
\texttt{lipschitz} & 108 & 317 & 108 & 317 & 108 & 317 & 108 & 317 \\
\texttt{swing\_cells} & 132 & 393 & 184 & 640 & 224 & 830 & 264 & 1020 \\
\texttt{swing\_direct} & 168 & 552 & 220 & 781 & 260 & 969 & 296 & 1175 \\
\midrule
\texttt{single\_bracket} & 1 & 1 & 1 & 1 & 1 & 1 & 1 & 1 \\
\texttt{fine\_grid} & \textbf{--} & \textbf{--} & \textbf{--} & \textbf{--} & \textbf{--} & \textbf{--} & \textbf{--} & \textbf{--} \\
\bottomrule
\end{tabular}
\end{center}

\paragraph{D:mult=3} 3 eigenvalues in $[478.7, 508.3]$; sweep took 542\,s.

\begin{center}
\small
\begin{tabular}{lrrrrrrrr}
\toprule
& \multicolumn{2}{c}{$\pr=10^{-4}$} & \multicolumn{2}{c}{$10^{-8}$} & \multicolumn{2}{c}{$10^{-11}$} & \multicolumn{2}{c}{$10^{-14}$} \\
\cmidrule(lr){2-3}\cmidrule(lr){4-5}\cmidrule(lr){6-7}\cmidrule(lr){8-9}
method & orc & aut & orc & aut & orc & aut & orc & aut \\
\midrule
\texttt{uniform\_grid} & 76 & 94 & 76 & 94 & 76 & 94 & 76 & 94 \\
\texttt{lappy} & 94 & 142 & 94 & 142 & 94 & 142 & 94 & 142 \\
\texttt{lappy\_repaired} & 221 & 221 & 331 & 1224 & 331 & 1224 & 331 & 1224 \\
\texttt{lipschitz} & 99 & 188 & 99 & 188 & 99 & 188 & 99 & 188 \\
\texttt{swing\_cells} & 119 & 261 & 158 & 352 & 188 & 422 & 218 & 492 \\
\texttt{swing\_direct} & 138 & 302 & 189 & 397 & 209 & 469 & 249 & 535 \\
\midrule
\texttt{single\_bracket} & 1 & 1 & 1 & 1 & 1 & 1 & 1 & 1 \\
\texttt{fine\_grid} & \textbf{--} & \textbf{--} & \textbf{--} & \textbf{--} & \textbf{--} & \textbf{--} & \textbf{--} & \textbf{--} \\
\bottomrule
\end{tabular}
\end{center}

\paragraph{D:mult=4} 4 eigenvalues in $[622.3, 660.8]$; sweep took 615\,s.

\begin{center}
\small
\begin{tabular}{lrrrrrrrr}
\toprule
& \multicolumn{2}{c}{$\pr=10^{-4}$} & \multicolumn{2}{c}{$10^{-8}$} & \multicolumn{2}{c}{$10^{-11}$} & \multicolumn{2}{c}{$10^{-14}$} \\
\cmidrule(lr){2-3}\cmidrule(lr){4-5}\cmidrule(lr){6-7}\cmidrule(lr){8-9}
method & orc & aut & orc & aut & orc & aut & orc & aut \\
\midrule
\texttt{uniform\_grid} & 76 & 94 & 76 & 94 & 76 & 94 & 76 & 94 \\
\texttt{lappy} & 94 & 142 & 94 & 142 & 94 & 142 & 94 & 142 \\
\texttt{lappy\_repaired} & 201 & 201 & 262 & 759 & 262 & 759 & 262 & 759 \\
\texttt{lipschitz} & 152 & 358 & 152 & 358 & 152 & 358 & 152 & 358 \\
\texttt{swing\_cells} & 146 & 454 & 185 & 597 & 215 & 707 & 245 & 817 \\
\texttt{swing\_direct} & 173 & 515 & 214 & 662 & 236 & 762 & 274 & 882 \\
\midrule
\texttt{single\_bracket} & 1 & 1 & 1 & 1 & 1 & 1 & 1 & 1 \\
\texttt{fine\_grid} & \textbf{--} & \textbf{--} & \textbf{--} & \textbf{--} & \textbf{--} & \textbf{--} & \textbf{--} & \textbf{--} \\
\bottomrule
\end{tabular}
\end{center}

\subsection{Family E --- mixed window}

\paragraph{E:mixed} 29 eigenvalues in $[150, 560]$; sweep took 1145\,s.

\begin{center}
\small
\begin{tabular}{lrrrrrrrr}
\toprule
& \multicolumn{2}{c}{$\pr=10^{-4}$} & \multicolumn{2}{c}{$10^{-8}$} & \multicolumn{2}{c}{$10^{-11}$} & \multicolumn{2}{c}{$10^{-14}$} \\
\cmidrule(lr){2-3}\cmidrule(lr){4-5}\cmidrule(lr){6-7}\cmidrule(lr){8-9}
method & orc & aut & orc & aut & orc & aut & orc & aut \\
\midrule
\texttt{uniform\_grid} & \textbf{--} & \textbf{--} & \textbf{--} & \textbf{--} & \textbf{--} & \textbf{--} & \textbf{--} & \textbf{--} \\
\texttt{lappy} & \textbf{--} & \textbf{--} & \textbf{--} & \textbf{--} & \textbf{--} & \textbf{--} & \textbf{--} & \textbf{--} \\
\texttt{lappy\_repaired} & \textbf{--} & \textbf{--} & \textbf{--} & \textbf{--} & \textbf{--} & \textbf{--} & \textbf{--} & \textbf{--} \\
\texttt{lipschitz} & \textbf{--} & \textbf{--} & \textbf{--} & \textbf{--} & \textbf{--} & \textbf{--} & \textbf{--} & \textbf{--} \\
\texttt{swing\_cells} & \textbf{--} & \textbf{--} & \textbf{--} & \textbf{--} & \textbf{--} & \textbf{--} & \textbf{--} & \textbf{--} \\
\texttt{swing\_direct} & \textbf{--} & \textbf{--} & \textbf{--} & \textbf{--} & \textbf{--} & \textbf{--} & \textbf{--} & \textbf{--} \\
\midrule
\texttt{single\_bracket} & \textbf{--} & \textbf{--} & \textbf{--} & \textbf{--} & \textbf{--} & \textbf{--} & \textbf{--} & \textbf{--} \\
\texttt{fine\_grid} & \textbf{--} & \textbf{--} & \textbf{--} & \textbf{--} & \textbf{--} & \textbf{--} & \textbf{--} & \textbf{--} \\
\bottomrule
\end{tabular}
\end{center}



\section{Analysis}

\subsection{The headline: nobody splits a pair below \texorpdfstring{$d/\lam \approx 10^{-4}$}{d/lam ~ 1e-4}}
\label{sec:headline}

Family B is the pair ladder. The table gives the \emph{oracle} cost --- the cheapest tested
setting that succeeded --- at each separation, at the tightest $\pr$ for which the pair must
still be \emph{split} rather than merged. A dash is failure at every setting within the
1500-evaluation budget.

\begin{center}
\begin{tabular}{rrrrrrrr}
\toprule
$d$ & $d/\lam$ & \texttt{uniform} & \texttt{lappy} & \texttt{lappy\_rep} & \texttt{lipschitz}
& \texttt{swing\_cells} & \texttt{swing\_direct} \\
\midrule
5.8e-01 & 1.2e-02 & 446 & 302 & 1316 & 293 & \textbf{131} & 152 \\
1.8e-01 & 3.6e-03 & --  & 1016 & 1103 & 1298 & \textbf{528} & 688 \\
5.9e-02 & 1.2e-03 & --  & -- & 922 & 1308 & \textbf{184} & 225 \\
1.8e-02 & 3.6e-04 & --  & -- & \textbf{830} & -- & -- & -- \\
5.9e-03 & 1.2e-04 & --  & -- & 1426 & -- & \textbf{658} & 844 \\
1.8e-03 & 3.6e-05 & --  & -- & -- & -- & -- & -- \\
\multicolumn{2}{r}{$\dots$ down to} & -- & -- & -- & -- & -- & -- \\
5.9e-12 & 1.2e-13 & --  & -- & -- & -- & -- & -- \\
1.8e-12 & 3.6e-14 & --  & -- & -- & -- & -- & -- \\
\bottomrule
\end{tabular}
\end{center}

\noindent
\textbf{Every method's splitting limit lies between $d/\lam = 1.2\mathrm{e}{-4}$ and
$3.6\mathrm{e}{-5}$} at this budget --- four to five orders coarser than the pencil's own
resolution limit of $\varepsilon/c \approx 4\mathrm{e}{-14}$ absolute, about $10^{-15}$
relative. \textbf{The gap between what the objective supports and what the searches extract is
roughly ten orders of magnitude}, and it is the most important number here.

Below the limit the methods do not fail loudly: they merge the pair and report success at any
$\pr$ loose enough to permit it. At $d = 1.8\mathrm{e}{-3}$ and $\pr = 10^{-4}$, where merging
is legal, \texttt{uniform\_grid} passes at 86 evaluations; at $\pr = 10^{-8}$ on the same case,
where it is not, every method fails.

\subsection{Family B is not the whole story, and reading it alone gives the wrong verdict}

A first draft of this report concluded from the ladder above that \texttt{lipschitz} should be
cut. The per-family breakdown contradicts that, and the contradiction is recorded here rather
than quietly fixed. Oracle\,/\,autonomous passes at $\pr = 10^{-8}$:

\begin{center}
\begin{tabular}{lrrrrr}
\toprule
method & A (1) & B (12) & C (8) & D (3) & E (1) \\
\midrule
\texttt{uniform\_grid}   & 1/0 & 3/2 & 3/2 & 3/3 & 0/0 \\
\texttt{lappy}           & 0/0 & 4/3 & 4/2 & 3/3 & 0/0 \\
\texttt{lappy\_repaired} & 1/0 & \textbf{6}/0 & \textbf{6}/2 & 3/3 & 0/0 \\
\texttt{lipschitz}       & 1/0 & 5/\textbf{5} & 5/\textbf{4} & 3/3 & 0/0 \\
\texttt{swing\_cells}    & 0/0 & \textbf{6}/\textbf{5} & 3/3 & 3/3 & 0/0 \\
\texttt{swing\_direct}   & 1/\textbf{1} & \textbf{6}/\textbf{5} & 3/3 & 3/3 & 0/0 \\
\bottomrule
\end{tabular}
\end{center}

Three things the pair ladder alone concealed:

\textbf{The swing methods are notably worse on clusters than on pairs} --- 3 of 8 on family C
against 6 of 12 on family B, where \texttt{lappy\_repaired} reaches 6 of 8 and
\texttt{lipschitz} 5. A cell decomposition cut at eigenvector crossings is well suited to
splitting two eigenvalues and less so to a tight group of three or four with unequal internal
gaps.

\textbf{\texttt{lappy\_repaired} has the best oracle record and the worst autonomous one} ---
6 of 12 on family B at oracle, and \textbf{0 of 12} autonomously. Essentially all of its
apparent capability is per-case tuning. That is the single largest oracle--autonomous gap in
the study, and it is exactly the failure mode the two-column protocol exists to expose: a
method that looks strongest in a tuned comparison and would ship as the weakest.

\textbf{\texttt{lipschitz} is the most consistent autonomous performer}, 5/5 on B and 4 of 5
on C --- its oracle and autonomous columns nearly coincide, meaning one fixed setting works
across cases. Its known theoretical weakness, that a pilot-estimated Lipschitz constant can
underestimate the true slope and render the completeness certificate false, did not bite on
rectangles. That is a real caveat for corner domains, not a measured failure here.

\textbf{Nobody passes family E} --- the 29-eigenvalue mixed window --- at any $\pr$ or any
setting. A window holding isolated levels, six pairs and a triple defeats every method in the
study within budget.

\subsection{Cost scales with \texorpdfstring{$\pr$}{prec} even when the requirement does not}

On most cases the purity requirement is \emph{identical} at $\pr = 10^{-8}$, $10^{-11}$ and
$10^{-14}$: the smallest true gap is far above $\pr\!\cdot\!\lam$ at all three, so nothing about
what must be resolved changes. Yet cost rises monotonically. On $d = 5.8\mathrm{e}{-1}$,
\texttt{swing\_cells} needs 131, 170, 200 and 230 evaluations at the four precisions --- a
$1.8\times$ increase for ten orders of $\pr$, buying nothing.

The cause is that every method bisects to $\pr\!\cdot\!\lam$ regardless of the local gap.
Refining a cell boundary to $10^{-13}$ when the nearest other eigenvalue is $0.58$ away is pure
waste. This violates the ``not more work than asked'' clause of the operational definition, and
it is a defect of the whole field rather than of any one method.

That this is cost and not capability was confirmed directly rather than inferred. On
$d = 5.9\mathrm{e}{-3}$ at $\pr = 10^{-14}$, where the requirement is identical to
$\pr = 10^{-11}$, \texttt{swing\_cells} fails at the 1500-evaluation cap and \textbf{succeeds at
1549 evaluations} once the cap is raised --- against the 1279 it needs at $10^{-11}$. A
21\,\% cost increase, and the only thing separating ``fails'' from ``succeeds'' is the budget
line.

\subsection{The controls validate the metrics}

\texttt{single\_bracket} returns $[(a,b)]$ and scores coverage $1.0$ with zero spurious on every
case. It is rejected by \textbf{purity alone}, on 18 of 21 cases. The three it passes are the
family-D exact-degeneracy cases, which contain a single level --- there one bracket \emph{is}
the correct answer, so the metric accepting it is right rather than broken.

\texttt{fine\_grid} exceeds the evaluation cap on every case and is rejected on \textbf{cost}.
Neither degenerate strategy survives, and each is rejected by the criterion intended for it.

\section{Filter verdict}

Capability claims within the measured budget of 1500 evaluations, on rectangles, for the
bracketing phase only.

\begin{description}
\item[\texttt{uniform\_grid} --- CUT as a candidate.] Weakest or joint-weakest in every family,
splits no pair tighter than $d/\lam = 1.2\mathrm{e}{-2}$. It should be retained as the control
that every proposal must beat, which is the role it already has.

\item[\texttt{lappy} (original) --- CUT.] Dominated by \texttt{lappy\_repaired} in every family
at both oracle and autonomous, and splits nothing below $d/\lam = 3.6\mathrm{e}{-3}$. Its
$\sig_2$-value recursion gate was already known to test the grid rather than the pair; this
measures the consequence.

\item[\texttt{lappy\_repaired} --- KEEP, with the gap as the priority.] Best oracle record of
any method (6/12 on B, 6/8 on C) and the worst autonomous (0/12 on B). Nothing about its
subdivision is wrong; what fails is its self-configuration. Closing that gap is worth more than
any change to its search strategy.

\item[\texttt{lipschitz} --- KEEP.] The most consistent autonomous performer, with oracle and
autonomous nearly coinciding. This reverses the draft verdict, which cut it on the pair ladder
alone. Its unsound-certificate risk is real but was not exercised here, and should be tested on
a corner domain before it is either trusted or discarded.

\item[\texttt{swing\_cells} --- KEEP.] Cheapest method wherever it succeeds, often by
$3$--$7\times$, and reaches the ladder's limit. Its weaknesses are clusters of three and four
(3 of 8) and the isolated-eigenvalue case, which it fails at every $\pr$.

\item[\texttt{swing\_direct} --- KEEP.] Costs 15--25\,\% more than \texttt{swing\_cells} and
buys no pair-splitting capability, but it is what rescues the isolated case (1/1 against 0/0).
The two mechanisms are orthogonal rather than ranked, and the ablation says so: cell-cutting
handles pairs, \textsc{direct} handles isolated wells.
\end{description}

\noindent
No method is cut on the strength of family B alone, and no method survives family E.

\section{What this cannot conclude}

\begin{itemize}
\item \textbf{A correction.} An earlier draft of this report claimed that
$\pr = 10^{-14}$ could never be exercised as a requirement on this instrument, on the grounds
that the only separation distinguishing it from $10^{-11}$ --- a pair at
$d/\lam \approx 10^{-13}$ --- sits just $1.2\times$ above the \texttt{snap} merge floor of the
closed-form spectrum. \textbf{That reasoning was wrong: family B does not use snapped truth.} It
takes the pair values from the mode formula directly, precisely so that a controlled separation
is not at the mercy of a merge policy. Measured, the two values are 833 floating-point ulps
apart at $\delta = 10^{-13}$ and 250 apart at $3 \times 10^{-14}$, and the pencil has
$d/(\varepsilon/c) = 138$ of headroom at the first.

Two such cases were therefore built and run, and both pass the floor precondition (headroom
$4.1$ and $3.3$). At $\pr = 10^{-4}$, $10^{-8}$ and $10^{-11}$ the pair may legally merge and
most methods pass cheaply; at $\pr = 10^{-14}$ it must be split, and \textbf{every method
fails}. So $\pr = 10^{-14}$ is now exercised as a genuine requirement, and the answer is
consistent with the ladder: nothing splits a pair at $d/\lam = 1.2\mathrm{e}{-13}$ when nothing
splits one at $3.6\mathrm{e}{-5}$.

\item \textbf{Rectangles only.} The property rectangles cannot reproduce is the tension floor
jumping between adjacent modes on a corner domain --- separately found to be largely an
artifact of a marginal basis, but re-measured at recipe grade on only one domain.

\item \textbf{Bracketing only.} A method that deliberately blurs bracketing into multiplicity
estimation is under-credited by construction. The swing is such a method: the same eigenvector
signal that cuts its cells also distinguishes a true degeneracy from two distinct eigenvalues,
and none of that is visible to a bracket-set metric.

\item \textbf{One budget, one seed, three settings per knob.} A method rejected here is
rejected \emph{at 1500 evaluations}, not in principle. The one non-monotonicity in the ladder
--- \texttt{lappy\_repaired} passing $d = 1.8\mathrm{e}{-2}$ while both swing variants pass the
tighter $5.9\mathrm{e}{-3}$ and fail it --- is most likely a knob-grid artifact, since
\texttt{per\_gap} was swept only over $\{6, 32, 160\}$, but that was not confirmed.

\item \textbf{The oracle column is best-of-$K$} and biased upward in the number of knobs. Only
the autonomous column may be compared across methods, and the gap between the two is the most
actionable quantity in the report.
\end{itemize}


\end{document}
